% TOP_PROBLEM: {"id": 3131, "title": "Are almost all graphs determined by their spectrum?", "category_id": 3, "category": {"id": 3, "name": "graph_theory", "display_name": "Graph Theory", "description": "Problems involving graphs, networks, and their properties.", "slug": "graph-theory", "order_index": 3, "created_at": "2026-07-31T15:26:25.671Z"}, "status": "open", "problem_number": "OPG-48368", "set_id": 12}
% FIRST_POSTED: 2026-10-01
% ATTEMPT_STATUS: unresolved
% UnsolvedMath ID 3131; source catalogue code OPG-48368
% !TeX program = lualatex
% Research attempt by GPT-6 Astra Ultra; no claim of a new complete solution.
\documentclass[11pt,a4paper]{article}
\usepackage[margin=25mm]{geometry}
\usepackage{fontspec}
\setmainfont{lmroman10-regular.otf}[BoldFont=lmroman10-bold.otf,
 ItalicFont=lmroman10-italic.otf,BoldItalicFont=lmroman10-bolditalic.otf]
\usepackage{amsmath,amssymb,mathtools}
\usepackage{unicode-math}
\setmathfont{latinmodern-math.otf}
\AtBeginDocument{\let\setminus\smallsetminus}
\usepackage{xurl,hyperref}
\hypersetup{colorlinks=true,urlcolor=blue}
\setcounter{secnumdepth}{0}
\setlength{\parindent}{0pt}
\setlength{\parskip}{5pt}
\setlength{\emergencystretch}{3em}
\begin{document}
\section*{Are almost all graphs determined by their spectrum?}\label{3131}
{\small UnsolvedMath ID 3131 · OPG-48368 · Graph Theory · OpenGarden}
\subsection{Definitions and mathematical statement}
For an integer $n\geq1$, let $\mathcal G_n$ be the set of all simple
undirected graphs on the labelled vertex set $[n]=\{1,\ldots,n\}$.
The adjacency matrix $A_G$ has entry $1$ at $(u,v)$ precisely when
$uv\in E(G)$, and has zero diagonal. Its adjacency spectrum is the
multiset of its $n$ real eigenvalues, including their multiplicities.
Two graphs are cospectral when these multisets agree; equivalently,
their characteristic polynomials $\det(tI-A_G)$ agree.

A graph $G$ is determined by its adjacency spectrum if every finite
simple graph $H$ cospectral with $G$ is isomorphic to $G$. Here an
isomorphism is a bijection between the vertex sets preserving adjacency
in both directions. Relabelling a graph therefore does not count as a
different spectral realization. Equality of spectra already forces
$|V(H)|=|V(G)|$.

The question is whether
\[
 \lim_{n\to\infty}
 \frac{|\{G\in\mathcal G_n:G\text{ is determined by its spectrum}\}|}
 {2^{\binom n2}}=1.
\]
Equivalently, choose each possible edge independently with probability
$1/2$; does the probability of a nonisomorphic cospectral mate tend to
zero? This specifies the usual labelled random-graph meaning of
``almost all.'' The data supplied are only the adjacency spectrum:
neither eigenvectors nor the spectrum of the complement or Laplacian
is included in the reconstruction requirement.
\subsection{Short English statement}
Does the adjacency spectrum uniquely identify a uniformly random large graph, up to renaming its vertices, with probability tending to one?
\subsection{Research attempt}
\paragraph{Scope and process.}
This is a GPT-6 Astra Ultra research attempt dated 1 October 2026, using
primary-source browsing and elementary mathematical checks. The canonical
definition above is retained. Results below are restricted results or
reductions, not a claim of a new solution to the entire catalogue question.
No formal proof assistant or external mathematical referee checked this text.


\paragraph{Literature and chosen obstruction.}
Koval and Kwan [R1] construct exponentially many graphs determined by
their adjacency spectrum, without claiming the almost-all conclusion.
Abiad, Brouwer and Haemers [R2] analyze Godsil--McKay switching and
the distinction between cospectrality and nonisomorphism. We check a
basic switching construction explicitly, then prove that one fixed-size
switching mechanism occurs with vanishing probability in $G(n,1/2)$.
Eliminating one mechanism is not a proof that all spectral mates vanish.

\paragraph{A matrix proof of switching.}
Let $X$ be a set of $k$ vertices, with $k$ even. Assume its induced
graph is regular and each vertex outside $X$ has either $0$, $k/2$
or $k$ neighbours in $X$. Write the adjacency matrix in blocks as
\[
 A=\begin{pmatrix}B&N\\N^{\mathsf T}&D\end{pmatrix},\qquad
 Q=\frac2k J-I_k.
\]
Since $J^2=kJ$, we have $Q^2=I_k$ and $Q^{\mathsf T}=Q$.
Regularity and symmetry of $B$ imply $BJ=JB=rJ$, whence
$QBQ=B$. A column $v$ of $N$ of weight zero or $k$ is fixed by
$Q$, and a column of weight $k/2$ is sent to ${\bf1}-v$.
Thus conjugation by $\operatorname{diag}(Q,I)$ replaces each
half-neighbourhood by its complement, leaves the other edges alone,
and yields a simple graph adjacency matrix $A'$. It is an orthogonal
similarity, so $A,A'$ have identical characteristic polynomials.

\paragraph{A verified nonisomorphic pair.}
Take independent $X=\{1,2,3,4\}$ and independent outside vertices
$a,b,c$, whose neighbourhoods in $X$ are respectively
$\{1,2\},\{1,3\},\{1,4\}$. The graph is a three-armed star with
each edge subdivided once. After switching, vertex $1$ is isolated
and the other six vertices form a cycle. The graphs are nonisomorphic
because their degree multisets are respectively
\[
 \{3,2,2,2,1,1,1\},\qquad\{2,2,2,2,2,2,0\}.
\]
The $3\times3$ matrix $N^{\mathsf T}N$ has diagonal entries two
and off-diagonal entries one, so equals $I_3+J_3$ and has eigenvalues
$4,1,1$. The bipartite adjacency spectrum therefore consists of
$\pm2$, two copies of each of $\pm1$, and one zero. Both
characteristic polynomials are
$t(t^2-4)(t^2-1)^2$. The exact rational-matrix check in
\texttt{scripts/check\_attempt\_batch\_algebra\_2026\_10\_01.py}
also verifies the conjugation and differing degree lists.

\paragraph{Why this local mechanism is rare in a random graph.}
Fix a particular four-vertex set $X$ in $G(n,1/2)$. For any outside
vertex, its four independent adjacency bits to $X$ have allowed
weights zero, two or four with probability
\[
 (1+6+1)/16=1/2.
\]
The bits for different outside vertices are independent. Therefore
the probability all outside vertices satisfy the switching condition
is $2^{-(n-4)}$. Requiring regularity inside $X$ can only decrease
this probability. A union bound gives
\[
 \mathbb P(\text{some four-vertex switching set exists})
 \le\binom n4\,2^{-(n-4)}\longrightarrow0.
\]
For a fixed even $k\ge4$, the same argument replaces $1/2$ by
$\theta_k=(2+\binom k{k/2})/2^k<1$, giving the bound
$\binom nk\theta_k^{n-k}\to0$. The case $k=2$ is different:
$Q$ just swaps the two vertices, so this operation produces an
isomorphic graph and is not a spectral ambiguity.

\paragraph{Verification and exact logical gap.}
The nonisomorphic example proves that spectra do not determine every
graph, but the conjecture concerns a probability tending to one.
The probability argument concerns each fixed size $k$; it cannot
be summed over unbounded $k$ without additional uniform estimates.
More fundamentally, a cospectral mate need not arise from a single
switching set of this form. Even ruling out all such switchings
would leave other orthogonal similarities and global constructions.
Adjacency cospectrality alone is used throughout; no eigenvectors,
Laplacian data or spectrum of the complement have been supplied.

\paragraph{Final status.}
\textbf{UNRESOLVED.} A genuine cospectral ambiguity and a rigorous
vanishing-probability estimate for each fixed-size switching mechanism
are established. Controlling every possible nonisomorphic cospectral
mate is the first missing universal step.

\subsection{Sources}
\begin{itemize}
\item[S1] ULAM.AI and the UnsolvedMath Contributors, supplied problems.json, record 3131 (OPG-48368), OpenGarden collection. Catalogue identity and source transcription; CC BY 4.0.
\url{https://huggingface.co/datasets/ulamai/UnsolvedMath}.
\item[S2] Open Problem Garden, Are almost all graphs determined by their spectrum?. Original problem and discussion; the mathematical formulation below is expanded from the supplied source record.
\url{http://www.openproblemgarden.org/op/are_almost_all_graphs_determined_by_their_spectrum}.

\item[R1] Illya Koval and Matthew Kwan, \emph{Exponentially many
graphs are determined by their spectrum}, arXiv:2309.09788v2,
2024 revision. Primary abstract checked 1 October 2026.
\url{https://arxiv.org/abs/2309.09788}.
\item[R2] Aida Abiad, Andries E. Brouwer and Willem H. Haemers,
\emph{Godsil--McKay switching and isomorphism}, arXiv:1406.4170,
2014. Primary research source for the switching framework.
\url{https://arxiv.org/abs/1406.4170}.

\end{itemize}
\end{document}
